% ==============================================================================
% T0 Theory / FFGFT: Document 183 (Chapter Content)
% Title: Google Willow and the Physics of Self-Measurement:
%        Topological Patterns, Quantum Echoes, and the Continuum of Complexity
% Author: Johann Pascher
% Date: April 2026
% Related Documents:
%   Doc. 159: Harmonic Structure of the T0 Torus
%   Doc. 164: T0 and 't Hooft CAI
%   Doc. 170: Discrete Thresholds (formulated there as working hypothesis)
%   Doc. 175: Qubit State Spaces
% External Sources:
%   Google Quantum AI / Nature (2024): Error correction below threshold
%   Will et al., Nature 645 (2025): Floquet topological order
%   Google Quantum AI / Nature (2025): Quantum Echoes / OTOC
% ==============================================================================

\section*{Abstract}

Google's Willow quantum processor achieved three experimental breakthroughs in 2024
and 2025 that carry particular significance from the perspective of the Fundamental
Fractal Geometric Field Theory (FFGFT). First, it was demonstrated for the first time
that quantum errors decrease exponentially as the code distance increases — a sign
that topologically stable field modes become experimentally accessible. Second,
topologically ordered states beyond thermal equilibrium were realized on Willow,
generating characteristic geometric patterns — chiral edge modes and anyons — that
cannot be simulated classically. Third, out-of-time-order correlators (Quantum Echoes)
revealed harmonic propagation of quantum information through constructive interference.
This document relates these findings to the FFGFT geometry of the 4D torus with
parameter $\xipar = 4/30{,}000$ and develops a position formulated in Document~170
as a working hypothesis: complexity and self-reference are not a binary on/off
phenomenon, but a physical continuum.

\clearpage

\section{The Three Willow Experiments}

\subsection{Error Correction Below Threshold (December 2024)}

In December 2024, Google Quantum AI published in \textit{Nature} the demonstration
that Willow crosses the threshold of quantum error correction for the first time
\cite{GoogleWillow2024}. The key result: the logical error rate decreases
exponentially with increasing code distance $d$. A $d=7$ code block achieves an
error rate that is $2.14$ times lower than a $d=5$ block.

Willow uses 105 superconducting transmon qubits in a square lattice, cooled to
approximately $10\,\text{mK}$. Physical single-qubit error rates are approximately
$0.05\,\%$ on Willow — compared to $0.10\,\%$ on the predecessor Sycamore
\cite{Will2025PubMed}. Error correction is based on a surface code, whose logical
qubits are encoded through topologically protected patterns.

\begin{keypoint}[Willow Breakthrough 1: Topological Stability]
More physical qubits reduce the logical error exponentially. This is the experimental
evidence that topologically stable configurations exist in quantum systems and are
accessible — not through external noise suppression, but through geometric order.
\end{keypoint}

\subsection{Floquet Topological Order: Patterns Beyond Equilibrium
(September 2025)}

Will et al.\ realised a \emph{Floquet topologically ordered phase} on Willow — a
quantum phase forbidden in thermal equilibrium but arising under periodic external
driving \cite{Will2025Nature}. This work appeared in \textit{Nature} in September 2025.

\begin{keyresult}[Measured Patterns on the Willow Processor]
The experiment imaged the following structures on a lattice of up to 58 qubits:
\begin{itemize}
    \item \textbf{Chiral edge modes:} Propagation of information along system
    boundaries in a preferred direction — a topological geometric signature that
    indicates classically irreproducible order.
    \item \textbf{Anyonic excitations:} Quasiparticles that \emph{transmute} into
    each other under Floquet driving — e-anyons become m-anyons and vice versa,
    a property of the non-abelian entanglement structure.
    \item \textbf{Topological invariant:} An interferometrically measured bulk
    invariant that quantifies the order and remains robust against
    integrability-breaking disorder.
\end{itemize}
\end{keyresult}

The authors write: \textit{``The non-equilibrium topological order that we have
probed in this work encapsulates a fundamentally unique phenomenology that is
forbidden in thermal equilibrium.''} \cite{Will2025Nature}
Classical matrix product state methods cannot efficiently simulate these
highly entangled dynamics; the quantum processor makes them directly observable.

Cochran et al.\ (2025) complemented this with a visualisation of charge and string
dynamics in (2+1)-dimensional lattice gauge theories on the same processor
\cite{Cochran2025Nature}.

\subsection{Quantum Echoes: Harmonic Propagation of Information (October 2025)}

In October 2025, Google Quantum AI demonstrated in \textit{Nature} the first
verifiable quantum advantage for a real algorithm: the \emph{Quantum Echoes}
(second-order out-of-time-order correlators, OTOC$^{(2)}$) \cite{GoogleEchoes2025}.

The algorithm works like a quantum echolocation system:
\begin{enumerate}
    \item \textbf{Forward evolution:} Entanglement of 65 qubits used.
    \item \textbf{Perturbation:} Local disturbance of a single ``butterfly'' qubit.
    \item \textbf{Time reversal:} Precise reverse evolution.
    \item \textbf{Measurement:} The echo is amplified through \textbf{constructive
    quantum interference}, yielding a measurable signal of information propagation.
\end{enumerate}

\begin{keypoint}[Quantum Echoes: Harmony as a Measured Signal]
The measurements show how quantum information propagates as a harmonic wave front
through the qubit lattice. The visualisations in \cite{GoogleEchoes2025} make this
propagation visible: concentric patterns spreading radially from the perturbation
qubit — a measurable harmonic pattern.
The OTOC$^{(2)}$ runs on Willow $13{,}000$ times faster than on the world's fastest
supercomputer Frontier (approximately 2.1 hours versus approximately 3.2 years).
\end{keypoint}

In a companion experiment with UC Berkeley, the algorithm was applied to compute the
geometry of organic molecules (15 and 28 atoms). The results agreed with NMR
measurements and provided additional structural information inaccessible to NMR
\cite{GoogleEchoes2025}.


\section{FFGFT Interpretation of the Three Experiments}

\subsection{The 4D Torus as Common Geometric Framework}

In FFGFT, spacetime is a 4D identification torus
$T^4 = \mathbb{R}^4/\mathbb{Z}^4$ with period $L_0 = \xipar \cdot \ell_P$.
The spectrum of the Laplace operator on this torus reads (Document~159):
\begin{equation}
    \lambda_{\mathbf{n}} = \frac{n_1^2 + n_2^2 + n_3^2 + n_4^2}{L_0^2},
    \qquad n_i \in \mathbb{Z}
\end{equation}
Only integer winding numbers $n_i$ generate stable modes. The sum over all modes
converges — not despite its infinitude, but because of the periodic geometry, in
analogy with Euler's Basel problem (Doc.~159):
\begin{equation}
    \zeta(2) = \sum_{n=1}^{\infty} \frac{1}{n^2} = \frac{\pi^2}{6}
\end{equation}

\textbf{Error correction} (Willow Breakthrough~1) is, in this language, the
experimental projection of physical states onto topologically stable torus modes.
The higher the code distance $d$, the more modes are engaged, and the more robust
the projection.

\textbf{Floquet topological order} (Breakthrough~2) corresponds to torus resonance
modes under periodic driving. A drive of frequency $\omega_F$ couples to the torus
geometry when
\begin{equation}
    \omega_F \cdot L_0 \approx \frac{n}{c}, \qquad n \in \mathbb{Z}
\end{equation}
i.e., when the driving frequency hits a torus mode. The resulting chiral edge modes
and anyons are geometrically forced, not accidental.

\textbf{Quantum Echoes} (Breakthrough~3) requires a more detailed treatment,
because it contains the deepest physical point.

\subsubsection*{``Instantaneous'' — Legitimate in Calculation, Misleading in Realisation}

Before analysing recursion and entanglement, a conceptual clarification is necessary
that is fundamental to understanding the Quantum Echoes experiment.

In quantum mechanics, there are two levels at which a system can be described:

\textbf{Level 1 — The mathematical description (timeless):}
The Hilbert space $\mathcal{H}$ of the 65-qubit system is a $2^{65}$-dimensional
complex vector space. In it, all possible states exist \emph{simultaneously as
mathematical objects} — independent of time. The energy eigenstates $|E_n\rangle$
of the Hamiltonian are timeless solutions of the stationary Schrödinger equation
$\hat{H}|E_n\rangle = E_n|E_n\rangle$.
When one writes in a calculation that all $2^{65}$ amplitudes are contained
``simultaneously'' in the wave function, this is \emph{legitimate}: it describes
the mathematical structure of the state, not a physical process. The wave function
\begin{equation}
    |\Psi\rangle = \sum_{k=1}^{2^{65}} c_k |k\rangle
\end{equation}
is an object in Hilbert space — timeless, complete, without sequence.

\textbf{Level 2 — Physical realisation (always in time):}
As soon as this mathematical object is realised on a physical system, time enters
inescapably. The time-dependent Schrödinger equation
\begin{equation}
    i\hbar \frac{\partial}{\partial t}|\Psi(t)\rangle = \hat{H}|\Psi(t)\rangle
\end{equation}
is a differential equation \emph{in time $t$}. Every time step of the Hamiltonian
requires physical time. Every quantum gate on Willow takes approximately 25--50
nanoseconds. A circuit of 23 steps and 65 qubits requires microseconds of coherence
time — and 2.1 hours of measurement time to collect sufficient statistical events.

\begin{important}[The Decisive Difference]
\textbf{In calculation}, ``instantaneous'' is legitimate: Hilbert space has no
inner time. All amplitudes exist simultaneously as a mathematical structure.
Statements such as ``all paths simultaneously in the amplitude'' are correct here.

\textbf{In realisation}, ``instantaneous'' is false and misleading: nothing happens
without time. Entanglement builds up step by step, bounded by the Lieb-Robinson
bound. The wave function evolves causally and locally — no quantum effect transmits
information faster than permitted. The 2.1 hours on Willow are \emph{real physical
time}, not a metaphor.
\end{important}

The advantage of the quantum computer lies precisely in this transition: it takes
the mathematical structure of Hilbert space — the timeless parallelism of all
amplitudes — and \emph{realises} it physically, step by step, in real time, through
unitary operators. The classical computer can only \emph{describe} this structure —
by computing the $2^{65}$ amplitudes one after another, which would take 3.2 years.
The quantum computer \emph{is} the structure — and that costs 2.1 hours.

No instantaneity. But genuine parallelism in the amplitude, physically realised
in real time.

\subsubsection*{Why Recursion and Mutual Influence Are Decisive}

In a classical system — a billiard table, say — when one ball is struck, the effect
propagates forward: ball A hits B, B hits C, and so on. The paths are independent;
they can be computed individually and summed. A classical computer follows exactly
this logic: it decomposes the problem into independent subproblems and adds the results.

A fully entangled quantum system works in a fundamentally different way.
No qubit is independent of the others — but this dependence does not arise
instantaneously. It builds up step by step, mediated by the interaction terms
in the Hamiltonian $\hat{H}$. The Schrödinger equation
\begin{equation}
    i\hbar \frac{\partial}{\partial t}|\Psi\rangle = \hat{H}|\Psi\rangle
\end{equation}
is a local, temporal differential equation. Entanglement between distant qubits
does not arise instantaneously but propagates at finite speed — bounded by the
Lieb-Robinson bound, which prescribes an effective information velocity in lattice
models such as Willow.

What quantum mechanics achieves is something different and more subtle: rather than
computing one path after another sequentially, the Hamiltonian
\emph{realises all recursion paths simultaneously in the total amplitude}. The wave
function $|\Psi(t)\rangle$ in the $2^{65}$-dimensional Hilbert space contains at
every moment the complete information about every possible sequence of qubit
interactions — weighted by the respective complex amplitude. This is not
instantaneity: it is \textbf{parallelism in the amplitude}, built up step by step
through unitary time evolution.

The decisive difference from a classical computer:
A classical algorithm must compute and store these $2^{65}$ amplitudes sequentially.
The quantum processor \emph{holds all of them simultaneously physically present} —
in the coherence of the quantum state — and propagates them through a single unitary
operator. The time this requires scales with the depth of the circuit, not with the
number of amplitudes.

\begin{technical}[Why This Is Classically Intractable]
A classical algorithm would need to track each of the $2^{65} \approx 10^{20}$
amplitudes exactly. For 65 qubits and 23 time steps, this amounts to approximately
$10^{20}$ complex numbers — more working memory than any existing supercomputer
possesses. Approximate methods (Monte Carlo, tensor networks) fail because the mutual
entanglement of all qubits exponentially amplifies approximation errors with system
size. This is not a technical but a principled problem: recursive interaction in a
fully entangled system is structurally incompressible into independent parts —
neither instantaneously nor sequentially.
\end{technical}

The echo arises because, after time reversal, the same recursion paths are traversed
in reverse. When the forward and backward evolution contains \emph{exactly the same
interactions in reversed order}, all amplitudes interfere constructively — the signal
emerges from the noise. Any deviation, any non-unitary interaction, destroys this
constructive interference and leads to destructive cancellation.

\begin{keypoint}[Quantum Echoes: Recursion Is the Measurement Device]
The quantum echo does not simply measure information propagation. It measures the
\textbf{depth and coherence of mutual influence}: how far back a system can
reconstruct its own past state when every element is simultaneously influenced by
all others.
The more complete the network of mutual coupling, the stronger the echo.
The greater the recursion depth, the more information is encoded in the echo — and
the more intractable the task becomes for any sequential algorithm.
\end{keypoint}

\subsubsection*{The FFGFT Interpretation: Self-Similarity of the Torus Spectrum}

In FFGFT geometry, this recursiveness is not an accidental product of the
experimental configuration. It is structural: the modes of the torus $T^4$ are
self-referential through the periodic boundary condition.
A torus mode $n$ has wavelength $\lambda_n = L_0/n$. The wavelength of mode $n$
is exactly the period of mode $n+1$ divided — the modes ``know'' each other through
their shared geometry.

The sum
\begin{equation}
    \zeta(2) = \sum_{n=1}^{\infty} \frac{1}{n^2} = \frac{\pi^2}{6}
\end{equation}
converges not despite infinitely many terms. It converges because each term $1/n^2$
\emph{weights} the influence of mode~$n$ by all lower modes: high modes contribute
little, low modes much — a natural hierarchy of mutual influence, inscribed in the
topology.

The harmonic wave front that the OTOC experiment makes visible — the concentric,
radially spreading interference patterns in the qubit lattice — is the experimental
image of this topological hierarchy. What one sees is not an arbitrary pattern: it is
the projection of the torus spectrum into the $10 \times 10$-qubit geometry of the
Willow chip.

The factor of $13{,}000$ over the classical supercomputer is therefore not an
efficiency advantage in the ordinary sense. It measures how far the parallelism
of amplitude propagation exceeds sequential computation. The quantum computer
takes 2.1 hours — it too requires time, it too follows the Schrödinger equation
step by step. But in each of those steps, a single unitary operator advances all
$2^{65}$ amplitudes simultaneously. The classical supercomputer would require
3.2 years because it must process these amplitudes sequentially.

\begin{foundation}[Physical Realisation vs.\ Simulation]
The decisive difference is not speed but structure.
The quantum computer \emph{realises} the recursion physically — through the unitary
time evolution of the Hamiltonian, which carries all entanglement paths simultaneously
in the amplitude.
The classical computer \emph{simulates} this amplitude sequentially — it can
describe the result, but it cannot be the object.
The factor of $13{,}000$ is the measurable expression of this structural difference,
measured in computation time on given hardware.
\end{foundation}


\section{Complexity as a Physical Continuum}

\subsection{Against the Levels Metaphor}

Document~170 presented a table of discrete complexity levels — from elementary
particles (Level~0) to consciousness (Level~4) — interpreting the fractal correction
$K_\text{frak}$ as a threshold operator. That presentation was formulated as a
\textit{working hypothesis} and included the caveat that the levels require further
formalisation. This document takes that caveat seriously and develops an alternative
perspective.

\begin{important}[Complexity as Continuum]
Complexity, self-reference, and what is called ``consciousness'' are not binary
phenomena with discrete thresholds in FFGFT.
They form a physical continuum that grows with the recursion depth of the
mass-time coupling $\tilde{T} \cdot m = 1$.
There is no sharp boundary beyond which a system ``has consciousness'' — any more
than there is a sharp boundary beyond which water is ``hot''.
\end{important}

\subsection{What the Willow Chip Illustrates on the Continuum}

The Willow chip is not a conscious system. But it is not structureless either:

\begin{itemize}
    \item \textbf{Single Josephson qubit:} The Josephson energy $E_J$ defines an
    effective mass and thereby a timescale $\tilde{T} = \hbar/E_J$. The qubit is a
    physical clock — the first step of self-reference.
    \item \textbf{Surface code lattice (105 qubits):} Collective interaction generates
    logical qubits that are robust against local perturbations. The system develops a
    form of ``memory'' — not cognitive, but physical: it preserves its own topological
    order against noise.
    \item \textbf{Floquet drive:} The system ``remembers'' the driving rhythm through
    time-crystalline structures. The drive period is inscribed in the topology of the
    system.
    \item \textbf{OTOC measurement:} The system measures its own information
    distribution. It references itself — geometrically precise, in measurable patterns.
\end{itemize}

This is a point on a continuum, not a threshold crossing. Between this point and the
human brain lies not an abyss but a continuous spectrum of increasing recursion depth,
integration breadth, and feedback stability.

\subsection{The Open Question — and the Honest Limit of Physics}

FFGFT describes this continuum geometrically: more qubits, deeper recursion, broader
entanglement move the system upward along the continuum. But physics cannot answer
where on this continuum what is called subjective experience arises — the experience
of colour, pain, time, meaning.

This is not a gap in the theory but an honest statement about the scope of physical
description. FFGFT can say precisely:

\begin{conclusionBox}[What FFGFT Can State]
\begin{itemize}
    \item The geometric structure $\xipar \cdot \ell_P$ that generates Willow patterns
    is the same structure from which all natural constants are derived.
    \item Complexity and self-reference grow continuously with the recursion depth of
    the mass-time coupling.
    \item The topological patterns measured in Willow experiments are not random
    artefacts but geometrically forced consequences of this structure.
\end{itemize}
What consciousness \textit{is} remains open — that is not a physical statement
expressible in an equation.
\end{conclusionBox}


\section{Summary}

\begin{table}[ht]
\centering
\begin{tabular}{|p{3.2cm}|p{4.8cm}|p{4.2cm}|}
\hline
\textbf{Experiment} & \textbf{Measured Finding} & \textbf{FFGFT Relation} \\
\hline
Error Correction\\Nature (Dec. 2024) &
Exponential error reduction\\with code distance &
Projection onto stable\\torus modes \\
\hline
Floquet topo.\ order\\Nature (Sep. 2025) &
Chiral edge modes,\\anyons, topo.\ invariant &
Torus resonance under\\Floquet drive \\
\hline
Quantum Echoes / OTOC\\Nature (Oct. 2025) &
Constructive interference,\\harmonic wave front &
Interferometry on torus;\\$\zeta(2) = \pi^2/6$ spectrum \\
\hline
\end{tabular}
\caption{Three Willow breakthroughs and their FFGFT interpretation}
\end{table}

The three experiments show consistently: quantum systems under controlled conditions
reveal geometric order — patterns that follow from the topology of the vacuum field.
This is the physically precise core behind intuitive descriptions associating quantum
computers with self-perception or patterns. The difference from speculative narrative:
the patterns are measured, quantified, and published in peer-reviewed journals.

\vspace{0.8cm}

\begin{thebibliography}{9}

\bibitem{GoogleWillow2024}
Google Quantum AI,
\textit{Quantum error correction below the surface code threshold},
\textit{Nature} (2024), DOI: 10.1038/s41586-024-08449-y

\bibitem{Will2025Nature}
M.~Will, T.\,A.~Cochran, E.~Rosenberg et al.,
\textit{Probing non-equilibrium topological order on a quantum processor},
\textit{Nature} \textbf{645}, 348--353 (2025),
DOI: 10.1038/s41586-025-09456-3

\bibitem{Will2025PubMed}
M.~Will et al.,
\textit{Probing non-equilibrium topological order on a quantum processor},
PubMed, PMID: 40931155 (2025)

\bibitem{Cochran2025Nature}
T.\,A.~Cochran, B.~Jobst, E.~Rosenberg et al.,
\textit{Visualizing dynamics of charges and strings in (2+1)D lattice gauge theories},
\textit{Nature} \textbf{642}, 315--320 (2025),
DOI: 10.1038/s41586-025-08999-9

\bibitem{GoogleEchoes2025}
Google Quantum AI and collaborators,
\textit{Observation of constructive interference at the edge of quantum ergodicity},
\textit{Nature} (2025), DOI: 10.1038/s41586-025-09526-6

\end{thebibliography}

\vspace{0.5cm}
\noindent
\textit{Draft for further elaboration. The connection between Floquet frequency
and torus fundamental frequency is a working hypothesis requiring formal verification.
Statements about the continuum of complexity are to be understood as a physically
motivated position, not a complete theory of consciousness.}
